\documentclass[12pt,a4paper]{report}

% =============================================================================
% KHAI BÁO CÁC PACKAGE CẦN THIẾT (ĐỘC LẬP, TƯƠNG THÍCH XELATEX & LUALATEX)
% =============================================================================
\usepackage{iftex}
\ifTUTeX
  \usepackage{fontspec}
  % Ưu tiên Times New Roman giống báo cáo gốc (trên Windows/macOS với XeLaTeX/LuaLaTeX),
  % tự động chuyển sang TeX Gyre Termes (bản chuẩn Times New Roman của TeX Live hỗ trợ đầy đủ tiếng Việt)
  \IfFontExistsTF{Times New Roman}{%
    \setmainfont{Times New Roman}%
  }{%
    \setmainfont{TeX Gyre Termes}%
  }
  \IfFontExistsTF{Arial}{%
    \setsansfont[Scale=0.92]{Arial}%
  }{%
    \setsansfont[Scale=0.92]{TeX Gyre Heros}%
  }
  \IfFontExistsTF{Consolas}{%
    \setmonofont[Scale=0.85]{Consolas}%
  }{%
    \setmonofont[Scale=0.85]{TeX Gyre Cursor}%
  }
\else
  \usepackage[utf8]{inputenc}
  \usepackage[T5]{fontenc}
  \usepackage{mathptmx}
\fi

\usepackage[a4paper,left=2.5cm,right=2.0cm,top=2.2cm,bottom=2.2cm,headheight=15pt]{geometry}
\usepackage{array}
\usepackage{longtable}
\usepackage{tabularx}
\usepackage{booktabs}
\usepackage[table]{xcolor}
\usepackage{colortbl}
\usepackage{caption}
\usepackage{fancyhdr}
\usepackage{hyperref}

% =============================================================================
% ĐỊNH NGHĨA MÀU SẮC & THIẾT LẬP ĐỊNG DẠNG THEO PHONG CÁCH BÁO CÁO GỐC
% =============================================================================
\definecolor{headerbg}{RGB}{230, 228, 248}
\definecolor{subhdrbg}{RGB}{242, 244, 248}
\definecolor{passgreen}{RGB}{18, 115, 55}
\definecolor{primaryblue}{RGB}{24, 62, 118}

\hypersetup{
  colorlinks=true,
  linkcolor=primaryblue,
  urlcolor=primaryblue,
  citecolor=primaryblue,
  pdftitle={Kiểm thử chức năng UC-01 và UC-02 - Hệ thống Quản lý Lớp học thêm},
  pdfauthor={Nhóm 01 - 2611COMP104405}
}

% Thiết lập nhãn tiếng Việt theo chuẩn báo cáo gốc
\renewcommand{\chaptername}{Chương}
\renewcommand{\contentsname}{Mục lục}
\renewcommand{\listtablename}{Danh mục bảng}
\renewcommand{\tablename}{Bảng}

% Giãn dòng và khoảng cách đoạn văn
\linespread{1.15}
\setlength{\parindent}{0pt}
\setlength{\parskip}{4pt plus 1pt minus 1pt}
\setlength{\LTpre}{4pt}
\setlength{\LTpost}{6pt}
\setlength{\tabcolsep}{4.5pt}
\renewcommand{\arraystretch}{1.14}
\arrayrulecolor{black}
\emergencystretch=3em

% Thiết lập caption bảng giống báo cáo gốc
\captionsetup[table]{
  labelfont=bf,
  labelsep=colon,
  justification=centering,
  singlelinecheck=true,
  skip=5pt
}

% =============================================================================
% CÁC LỆNH HỖ TRỢ TRÌNH BÀY BẢNG KIỂM THỬ
% =============================================================================
\newcolumntype{L}[1]{>{\raggedright\setlength{\parskip}{1pt}\emergencystretch=2em\arraybackslash}p{#1}}
\newcolumntype{C}[1]{>{\centering\setlength{\parskip}{1pt}\arraybackslash}p{#1}}
\newcolumntype{R}[1]{>{\raggedleft\setlength{\parskip}{1pt}\arraybackslash}p{#1}}

% Lệnh hiển thị trạng thái PASS nổi bật, rõ ràng
\newcommand{\passbadge}{\textbf{\textcolor{passgreen}{PASS}}}

% Lệnh định dạng khối thông tin mô tả chi tiết trong cột 2 của bảng kiểm thử
\newcommand{\tcheader}[2]{%
  \textbf{#1} \textit{(Mã: #2)}\par
}

\newcommand{\tcmeta}[2]{%
  \textbf{Chức năng:} #1 -- \textbf{Vai trò:} #2\par
}

\newcommand{\tcpre}[1]{%
  \textbf{ĐK (Tiền điều kiện):} #1\par
}

\newcommand{\tcdata}[1]{%
  \textbf{Dữ liệu kiểm thử:} #1\par
}

\begin{document}

% Thiết lập bộ đếm chương = 6 để khớp trực tiếp với Chương 6 của báo cáo gốc
\setcounter{chapter}{5}
\chapter{KIỂM THỬ CHƯƠNG TRÌNH}
%=============================================================================
% MỤC 6.1: KIỂM THỬ CÁC CHỨC NĂNG ĐÃ THỰC HIỆN
% =============================================================================
\section{Kiểm thử các chức năng đã thực hiện}

% -----------------------------------------------------------------------------
% PHẦN 1: KIỂM THỬ CHO CHỨC NĂNG 1 (UC-01 - ĐĂNG NHẬP)
% -----------------------------------------------------------------------------
\subsection{Kiểm thử cho chức năng 1 - Đăng nhập (UC-01)}
\label{subsec:test_uc01}

Chức năng Đăng nhập (\textbf{UC-01}) cho phép Giáo viên (\texttt{ACT-01}), Học sinh (\texttt{ACT-02}) và Phụ huynh (\texttt{ACT-03}) xác thực danh tính bằng tài khoản đã được cấp phát (hệ thống không mở đăng ký công khai). Sau khi đối chiếu hợp lệ qua lớp \texttt{DieuKhienXacThuc} và bảng \texttt{users}, hệ thống khởi tạo phiên làm việc (\texttt{Session}) và điều hướng người dùng đến giao diện Tổng quan (\textit{Dashboard}) tương ứng với vai trò.

{\small\linespread{1.05}\selectfont
\begin{longtable}{|C{0.8cm}|L{6.8cm}|L{3.0cm}|L{2.9cm}|C{1.2cm}|}
\caption{\normalsize Kết quả kiểm thử chức năng đăng nhập (UC-01)}
\label{tab:test_uc01} \\
\hline
\rowcolor{headerbg}
\textbf{STT} & \textbf{Tên Test Case \& Mô tả} & \textbf{Đầu ra (Output) mong đợi} & \textbf{Ngoại lệ / Thực tế} & \textbf{Kết quả} \\
\hline
\endfirsthead

\multicolumn{5}{l}{\textit{  } \\[3pt]
\hline
\rowcolor{headerbg}
\textbf{STT} & \textbf{Tên Test Case \& Mô tả} & \textbf{Đầu ra (Output) mong đợi} & \textbf{Ngoại lệ / Thực tế} & \textbf{Kết quả} \\
\hline
\endhead

\hline
\multicolumn{5}{r}{\textit{ }} \\
\endfoot

\hline
\endlastfoot

% --- TC_UC01_001 ---
1 &
\tcheader{Đăng nhập thành công bằng tài khoản Giáo viên}{TC\_UC01\_001}
\tcmeta{UC-01 (Đăng nhập)}{Giáo viên (ACT-01)}
\tcpre{Hệ thống hoạt động bình thường; tài khoản Giáo viên đã tồn tại và ở trạng thái hoạt động (\texttt{is\_active = true}).}
\tcdata{Tên đăng nhập: \texttt{giaovien@\allowbreak lopminh.edu.vn}; Mật khẩu: \texttt{123}.}
\textbf{HĐ (Các bước thực hiện):}\par
1. Mở trình duyệt và truy cập trang Đăng nhập của hệ thống.\par
2. Nhập tên đăng nhập hợp lệ \texttt{giaovien@\allowbreak lopminh.edu.vn} vào ô \textit{Tên đăng nhập hoặc email}.\par
3. Nhập mật khẩu hợp lệ \texttt{123} vào ô \textit{Mật khẩu}.\par
4. Nhấn nút \textbf{Đăng nhập}.\par
5. Kiểm tra phiên làm việc và giao diện hiển thị sau khi đăng nhập.
&
Hệ thống xác thực thành công, khởi tạo phiên làm việc với vai trò \texttt{TEACHER} và chuyển hướng đến giao diện \textbf{Tổng quan của không gian giáo viên} (hiển thị lời chào giáo viên, các thẻ thống kê và thanh điều hướng quản lý lớp học, tài khoản, điểm danh, điểm số, học phí).
&
Không có ngoại lệ. Hệ thống xác thực đúng tài khoản Giáo viên, tạo phiên đăng nhập hợp lệ và hiển thị chính xác giao diện Tổng quan dành cho Giáo viên.
&
\passbadge \\
\hline

% --- TC_UC01_002 ---
2 &
\tcheader{Đăng nhập thành công bằng tài khoản Học sinh}{TC\_UC01\_002}
\tcmeta{UC-01 (Đăng nhập)}{Học sinh (ACT-02)}
\tcpre{Tài khoản Học sinh đã được Giáo viên khởi tạo, cấp phát và đang ở trạng thái hoạt động (\texttt{is\_active = true}).}
\tcdata{Tên đăng nhập: \texttt{hocsinh@\allowbreak lopminh.edu.vn}; Mật khẩu: \texttt{123}.}
\textbf{HĐ (Các bước thực hiện):}\par
1. Truy cập trang Đăng nhập của hệ thống.\par
2. Nhập tên đăng nhập hợp lệ của Học sinh là \texttt{hocsinh@\allowbreak lopminh.edu.vn} vào ô \textit{Tên đăng nhập hoặc email}.\par
3. Nhập mật khẩu hợp lệ \texttt{123} vào ô \textit{Mật khẩu}.\par
4. Nhấn nút \textbf{Đăng nhập}.\par
5. Quan sát màn hình trang chủ và thanh điều hướng được tải lên.
&
Hệ thống xác thực thành công, cấp phiên làm việc với vai trò \texttt{STUDENT} và chuyển hướng đến giao diện \textbf{Tổng quan của không gian học sinh} (hiển thị các mục: Tổng quan, Lịch học, Điểm danh, Kết quả học tập, Bài tập, Thông báo, Hồ sơ cá nhân).
&
Không có ngoại lệ. Hệ thống đăng nhập thành công và điều hướng đúng về cổng thông tin học tập cá nhân dành cho Học sinh.
&
\passbadge \\
\hline

% --- TC_UC01_003 ---
3 &
\tcheader{Đăng nhập thành công bằng tài khoản Phụ huynh}{TC\_UC01\_003}
\tcmeta{UC-01 (Đăng nhập)}{Phụ huynh (ACT-03)}
\tcpre{Tài khoản Phụ huynh đã được Giáo viên khởi tạo, liên kết với hồ sơ Học sinh và đang hoạt động (\texttt{is\_active = true}).}
\tcdata{Tên đăng nhập: \texttt{phuhuynh@\allowbreak lopminh.edu.vn}; Mật khẩu: \texttt{123}.}
\textbf{HĐ (Các bước thực hiện):}\par
1. Truy cập trang Đăng nhập của hệ thống.\par
2. Nhập tên đăng nhập hợp lệ của Phụ huynh là \texttt{phuhuynh@\allowbreak lopminh.edu.vn} vào ô \textit{Tên đăng nhập hoặc email}.\par
3. Nhập mật khẩu hợp lệ \texttt{123} vào ô \textit{Mật khẩu}.\par
4. Nhấn nút \textbf{Đăng nhập}.\par
5. Kiểm tra giao diện trang chủ và danh mục chức năng hiển thị.
&
Hệ thống xác thực thành công, khởi tạo phiên làm việc với vai trò \texttt{PARENT} và chuyển hướng đến giao diện \textbf{Tổng quan của không gian phụ huynh} (hiển thị thông tin học sinh liên kết, chuyên cần, kết quả học tập và học phí).
&
Không có ngoại lệ. Hệ thống xác thực đúng vai trò Phụ huynh và hiển thị chính xác giao diện Tổng quan dành cho Phụ huynh.
&
\passbadge \\
\hline

% --- TC_UC01_004 ---
4 &
\tcheader{Từ chối đăng nhập khi sai mật khẩu}{TC\_UC01\_004}
\tcmeta{UC-01 (Đăng nhập)}{Mọi vai trò}
\tcpre{Tài khoản \texttt{giaovien@\allowbreak lopminh.edu.vn} đã tồn tại trên hệ thống với mật khẩu đúng là \texttt{123}.}
\tcdata{Tên đăng nhập: \texttt{giaovien@\allowbreak lopminh.edu.vn}; Mật khẩu nhập vào: \texttt{SaiMatKhau999}.}
\textbf{HĐ (Các bước thực hiện):}\par
1. Truy cập trang Đăng nhập của hệ thống.\par
2. Nhập tên đăng nhập đúng \texttt{giaovien@\allowbreak lopminh.edu.vn} vào ô \textit{Tên đăng nhập hoặc email}.\par
3. Nhập mật khẩu không chính xác \texttt{SaiMatKhau999} vào ô \textit{Mật khẩu}.\par
4. Nhấn nút \textbf{Đăng nhập}.\par
5. Kiểm tra thông báo phản hồi trên màn hình và trạng thái phiên làm việc.
&
Hệ thống từ chối yêu cầu đăng nhập, giữ người dùng ở lại trang Đăng nhập, không khởi tạo phiên làm việc và hiển thị thông báo lỗi: \textit{``Tài khoản hoặc mật khẩu không chính xác''}.
&
Hệ thống đối chiếu mã băm mật khẩu không khớp, từ chối đăng nhập, hiển thị thông báo lỗi \textit{``Tài khoản hoặc mật khẩu không chính xác''} và không tạo phiên làm việc.
&
\passbadge \\
\hline

% --- TC_UC01_005 ---
5 &
\tcheader{Từ chối đăng nhập khi sai tên đăng nhập}{TC\_UC01\_005}
\tcmeta{UC-01 (Đăng nhập)}{Người dùng}
\tcpre{Trang Đăng nhập sẵn sàng tiếp nhận dữ liệu; tên đăng nhập \texttt{khongtontai@\allowbreak lopminh.edu.vn} chưa từng được tạo trong bảng \texttt{users}.}
\tcdata{Tên đăng nhập: \texttt{khongtontai@\allowbreak lopminh.edu.vn}; Mật khẩu: \texttt{123}.}
\textbf{HĐ (Các bước thực hiện):}\par
1. Truy cập trang Đăng nhập của hệ thống.\par
2. Nhập chuỗi \texttt{khongtontai@\allowbreak lopminh.edu.vn} vào ô \textit{Tên đăng nhập hoặc email}.\par
3. Nhập mật khẩu \texttt{123} vào ô \textit{Mật khẩu}.\par
4. Nhấn nút \textbf{Đăng nhập}.\par
5. Kiểm tra phản hồi của hệ thống trên giao diện.
&
Hệ thống không cho phép đăng nhập, không tiết lộ việc tên tài khoản có tồn tại riêng lẻ hay không mà hiển thị thông báo chuẩn: \textit{``Tài khoản hoặc mật khẩu không chính xác''} và yêu cầu nhập lại.
&
Hệ thống từ chối đăng nhập, hiển thị đúng thông báo lỗi \textit{``Tài khoản hoặc mật khẩu không chính xác''} và giữ nguyên màn hình Đăng nhập.
&
\passbadge \\
\hline

% --- TC_UC01_006 ---
6 &
\tcheader{Kiểm tra dữ liệu đăng nhập bỏ trống}{TC\_UC01\_006}
\tcmeta{UC-01 (Đăng nhập)}{Người dùng}
\tcpre{Người dùng đang mở màn hình Đăng nhập tài khoản; các ô nhập liệu đều đang trống.}
\tcdata{Tên đăng nhập: \textit{(Để trống -- \texttt{""})}; Mật khẩu: \textit{(Để trống -- \texttt{""})}.}
\textbf{HĐ (Các bước thực hiện):}\par
1. Truy cập trang Đăng nhập của hệ thống.\par
2. Để trống hoàn toàn ô \textit{Tên đăng nhập hoặc email}.\par
3. Để trống hoàn toàn ô \textit{Mật khẩu}.\par
4. Nhấn trực tiếp nút \textbf{Đăng nhập}.\par
5. Quan sát trạng thái kiểm tra hợp lệ của biểu mẫu.
&
Hệ thống ngăn gửi yêu cầu xác thực rỗng lên máy chủ, đánh dấu các trường bắt buộc và hiển thị thông báo yêu cầu người dùng nhập tên đăng nhập và mật khẩu trước khi tiếp tục.
&
Hệ thống chặn thao tác gửi biểu mẫu rỗng ngay tại giao diện, hiển thị cảnh báo yêu cầu nhập thông tin ở cả hai trường bắt buộc và không tạo phiên đăng nhập.
&
\passbadge \\
\hline

% --- TC_UC01_007 ---
7 &
\tcheader{Kiểm tra một trường đăng nhập bị bỏ trống}{TC\_UC01\_007}
\tcmeta{UC-01 (Đăng nhập)}{Người dùng}
\tcpre{Người dùng đang ở màn hình Đăng nhập tài khoản.}
\tcdata{TH1: Tên đăng nhập = \texttt{giaovien@\allowbreak lopminh.edu.vn}, Mật khẩu = \textit{(Trống)}; TH2: Tên đăng nhập = \textit{(Trống)}, Mật khẩu = \texttt{123}.}
\textbf{HĐ (Các bước thực hiện):}\par
1. Truy cập trang Đăng nhập của hệ thống.\par
2. Nhập \texttt{giaovien@\allowbreak lopminh.edu.vn} vào ô \textit{Tên đăng nhập hoặc email}, để trống ô \textit{Mật khẩu} và nhấn \textbf{Đăng nhập}.\par
3. Quan sát thông báo lỗi tại trường \textit{Mật khẩu}.\par
4. Xóa ô \textit{Tên đăng nhập hoặc email}, chỉ nhập \texttt{123} vào ô \textit{Mật khẩu} và nhấn \textbf{Đăng nhập}.\par
5. Kiểm tra thông báo phản hồi tại trường bị bỏ trống.
&
Trong cả hai trường hợp, hệ thống đều không thực hiện đăng nhập và hiển thị thông báo yêu cầu điền đầy đủ trường thông tin còn thiếu (Mật khẩu hoặc Tên đăng nhập).
&
Hệ thống kiểm tra ràng buộc dữ liệu chính xác, báo lỗi ngay tại trường bị bỏ trống trong cả hai trường hợp và không cho phép đăng nhập.
&
\passbadge \\
\hline

% --- TC_UC01_008 ---
8 &
\tcheader{Ngăn truy cập trang nội bộ khi chưa đăng nhập}{TC\_UC01\_008}
\tcmeta{UC-01 (Đăng nhập)}{Khách (Chưa đăng nhập)}
\tcpre{Trình duyệt chưa có phiên làm việc (\texttt{Session Token} trống hoặc đã đăng xuất khỏi hệ thống).}
\tcdata{URL nội bộ thử truy cập trực tiếp: \texttt{/pages/teacher/\allowbreak dashboard.html}, \texttt{/pages/teacher/\allowbreak accounts.html}.}
\textbf{HĐ (Các bước thực hiện):}\par
1. Đảm bảo người dùng đã đăng xuất hoàn toàn khỏi hệ thống (xóa phiên làm việc hiện tại).\par
2. Nhập trực tiếp đường dẫn trang nội bộ (ví dụ: trang Tổng quan hoặc trang Quản lý tài khoản của Giáo viên) lên thanh địa chỉ trình duyệt và nhấn \texttt{Enter}.\par
3. Kiểm tra cơ chế chặn của \texttt{Auth Guard} và màn hình đích được hiển thị.
&
Hệ thống từ chối quyền truy cập vào trang nội bộ khi chưa xác thực, chặn hiển thị dữ liệu nghiệp vụ và tự động chuyển hướng trình duyệt về trang \textbf{Đăng nhập tài khoản}.
&
Bộ kiểm soát \texttt{Auth Guard} và \texttt{Session Middleware} chặn truy cập trái phép và điều hướng người dùng về trang Đăng nhập theo đúng thiết kế bảo mật.
&
\passbadge \\
\hline
\end{longtable}
}

\textbf{Bảng tổng kết kiểm thử chức năng UC-01:}

Kết quả thực hiện kiểm thử đối với chức năng \textbf{UC-01: Đăng nhập} gồm 08 kịch bản kiểm thử được tổng hợp tại Bảng~\ref{tab:summary_uc01}.

{\small
\begin{longtable}{|C{2.9cm}|C{2.8cm}|C{2.8cm}|C{2.8cm}|C{2.8cm}|}
\caption{\normalsize Bảng tổng kết kết quả kiểm thử chức năng Đăng nhập (UC-01)}
\label{tab:summary_uc01} \\
\hline
\rowcolor{headerbg}
\textbf{Tổng số test case} & \textbf{Số test case PASS} & \textbf{Số test case FAIL} & \textbf{Tỷ lệ PASS} & \textbf{Tỷ lệ FAIL} \\
\hline
\endfirsthead
\hline
\endfoot
8 & \textbf{\textcolor{passgreen}{8}} & 0 & \textbf{\textcolor{passgreen}{100\%}} & 0\% \\
\hline
\end{longtable}
}

% -----------------------------------------------------------------------------
% PHẦN 2: KIỂM THỬ CHO CHỨC NĂNG 2 (UC-02 - QUẢN LÝ NGƯỜI DÙNG VÀ PHÂN QUYỀN)
% -----------------------------------------------------------------------------
\subsection{Kiểm thử cho chức năng 2 - Quản lý người dùng và phân quyền (UC-02)}
\label{subsec:test_uc02}

Chức năng Quản lý người dùng và phân quyền (\textbf{UC-02}) hỗ trợ Giáo viên (\texttt{ACT-01}) khởi tạo, cập nhật, xóa/vô hiệu hóa tài khoản Học sinh, Phụ huynh và thiết lập liên kết giám hộ qua bảng \texttt{parent\_\allowbreak student\_\allowbreak links}. Cơ chế phân quyền (\texttt{TEACHER}, \texttt{STUDENT}, \texttt{PARENT}) được kiểm soát đồng bộ ở cả giao diện và tầng nghiệp vụ, đảm bảo chỉ Giáo viên mới có quyền thao tác quản trị người dùng.

{\small\linespread{1.05}\selectfont
\begin{longtable}{|C{0.8cm}|L{6.8cm}|L{3.0cm}|L{2.9cm}|C{1.2cm}|}
\caption{\normalsize Kết quả kiểm thử chức năng quản lý người dùng và phân quyền (UC-02)}
\label{tab:test_uc02} \\
\hline
\rowcolor{headerbg}
\textbf{STT} & \textbf{Tên Test Case \& Mô tả} & \textbf{Đầu ra (Output) mong đợi} & \textbf{Ngoại lệ / Thực tế} & \textbf{Kết quả} \\
\hline
\endfirsthead

\multicolumn{5}{l}{\textit{  } \\[3pt]
\hline
\rowcolor{headerbg}
\textbf{STT} & \textbf{Tên Test Case \& Mô tả} & \textbf{Đầu ra (Output) mong đợi} & \textbf{Ngoại lệ / Thực tế} & \textbf{Kết quả} \\
\hline
\endhead

\hline
\multicolumn{5}{r}{\textit{ }} \\
\endfoot

\hline
\endlastfoot

% --- TC_UC02_001 ---
1 &
\tcheader{Tạo tài khoản Học sinh hợp lệ}{TC\_UC02\_001}
\tcmeta{UC-02 (QL người dùng)}{Giáo viên (ACT-01)}
\tcpre{Giáo viên đã đăng nhập (\texttt{giaovien@\allowbreak lopminh.edu.vn}); lớp \textit{Toán 10 $\cdot$ Nền tảng (TO10A)} đang hoạt động; tài khoản \texttt{hocsinh@\allowbreak lopminh.edu.vn} chưa tồn tại.}
\tcdata{Vai trò: \texttt{Học sinh}; Họ tên: \texttt{Nguyen Van A}; Ngày sinh: \texttt{15/05/2010}; SĐT: \texttt{0912345678}; Tên ĐN/Email: \texttt{hocsinh@\allowbreak lopminh.edu.vn}; Mật khẩu: \texttt{123}.}
\textbf{HĐ (Các bước thực hiện):}\par
1. Đăng nhập bằng tài khoản Giáo viên (\texttt{giaovien@\allowbreak lopminh.edu.vn} / \texttt{123}) và chọn mục \textbf{Tài khoản} trên thanh điều hướng.\par
2. Nhấn nút \textbf{Tạo tài khoản} ở góc trên bên phải màn hình.\par
3. Chọn vai trò \textbf{Học sinh}, nhập đầy đủ Họ và tên, Ngày sinh, Số điện thoại và chọn lớp học tham gia.\par
4. Nhập Tên đăng nhập \texttt{hocsinh@\allowbreak lopminh.edu.vn} và Mật khẩu tạm thời \texttt{123}.\par
5. Nhấn nút \textbf{Lưu / Tạo tài khoản} và kiểm tra danh sách tài khoản.
&
Hệ thống khởi tạo thành công tài khoản Học sinh mới với vai trò \texttt{STUDENT}, lưu mật khẩu dạng băm vào bảng \texttt{users}, khởi tạo hồ sơ học sinh tương ứng và hiển thị tài khoản \texttt{Nguyen Van A} ở trạng thái \textit{Hoạt động} trên danh sách.
&
Không có ngoại lệ. Tài khoản \texttt{hocsinh@\allowbreak lopminh.edu.vn} được tạo thành công, hiển thị đầy đủ trên danh sách tài khoản và có thể dùng để đăng nhập.
&
\passbadge \\
\hline

% --- TC_UC02_002 ---
2 &
\tcheader{Tạo tài khoản Phụ huynh và liên kết Học sinh}{TC\_UC02\_002}
\tcmeta{UC-02 (QL người dùng)}{Giáo viên (ACT-01)}
\tcpre{Giáo viên đã đăng nhập; hồ sơ Học sinh \texttt{Nguyen Van A} (mã \texttt{HS001}, tài khoản \texttt{hocsinh@\allowbreak lopminh.edu.vn}) đã tồn tại trên hệ thống.}
\tcdata{Vai trò: \texttt{Phụ huynh}; Họ tên: \texttt{Nguyen Van B}; Quan hệ: \texttt{Bố}; SĐT: \texttt{0903456789}; Tên ĐN/Email: \texttt{phuhuynh@\allowbreak lopminh.edu.vn}; Mật khẩu: \texttt{123}; Liên kết HS: \texttt{Nguyen Van A (HS001)}.}
\textbf{HĐ (Các bước thực hiện):}\par
1. Tại màn hình \textbf{Quản lý tài khoản}, nhấn nút \textbf{Tạo tài khoản}.\par
2. Chọn vai trò \textbf{Phụ huynh}, nhập Họ và tên \texttt{Nguyen Van B}, Số điện thoại, Tên đăng nhập \texttt{phuhuynh@\allowbreak lopminh.edu.vn} và Mật khẩu \texttt{123}.\par
3. Tại mục \textbf{Liên kết học sinh}, tìm và chọn học sinh \texttt{Nguyen Van A (HS001)} cùng mối quan hệ là \texttt{Bố}.\par
4. Nhấn nút \textbf{Lưu} để hoàn tất khởi tạo.\par
5. Kiểm tra cột \textit{Liên kết gia đình} trên bảng danh sách tài khoản.
&
Hệ thống khởi tạo tài khoản Phụ huynh với vai trò \texttt{PARENT} trong bảng \texttt{users}, đồng thời ghi nhận bản ghi liên kết vào bảng \texttt{parent\_\allowbreak student\_\allowbreak links} gắn với học sinh \texttt{HS001} và hiển thị đúng tên học sinh liên kết trên giao diện.
&
Không có ngoại lệ. Hệ thống tạo tài khoản \texttt{phuhuynh@\allowbreak lopminh.edu.vn} thành công và thiết lập chính xác liên kết gia đình với học sinh \texttt{Nguyen Van A}.
&
\passbadge \\
\hline

% --- TC_UC02_003 ---
3 &
\tcheader{Không cho lưu tài khoản Phụ huynh khi chưa liên kết Học sinh}{TC\_UC02\_003}
\tcmeta{UC-02 (QL người dùng)}{Giáo viên (ACT-01)}
\tcpre{Giáo viên đã đăng nhập và đang mở biểu mẫu tạo mới tài khoản người dùng.}
\tcdata{Vai trò: \texttt{Phụ huynh}; Họ tên: \texttt{Tran Thi C}; SĐT: \texttt{0908112233}; Tên ĐN: \texttt{phuhuynh@\allowbreak lopminh.edu.vn}; Mật khẩu: \texttt{123}; Học sinh liên kết: \textit{(Chưa chọn -- Bỏ trống)}.}
\textbf{HĐ (Các bước thực hiện):}\par
1. Truy cập giao diện \textbf{Quản lý tài khoản} và nhấn \textbf{Tạo tài khoản}.\par
2. Chọn vai trò là \textbf{Phụ huynh}.\par
3. Nhập đầy đủ Họ và tên, Số điện thoại, Tên đăng nhập \texttt{phuhuynh@\allowbreak lopminh.edu.vn} và Mật khẩu \texttt{123}.\par
4. Không chọn bất kỳ học sinh nào tại khung \textbf{Liên kết học sinh}.\par
5. Nhấn nút \textbf{Lưu} và kiểm tra thông báo ràng buộc của hệ thống.
&
Hệ thống ngăn không cho lưu tài khoản Phụ huynh khi chưa có học sinh liên kết, hiển thị thông báo yêu cầu Giáo viên chọn ít nhất một tài khoản Học sinh để liên kết theo đúng quy trình \texttt{AD-02}.
&
Hệ thống chặn thao tác lưu, hiển thị cảnh báo bắt buộc chọn Học sinh cần liên kết và không sinh bản ghi Phụ huynh mồ côi trong cơ sở dữ liệu.
&
\passbadge \\
\hline

% --- TC_UC02_004 ---
4 &
\tcheader{Kiểm tra thông tin bắt buộc khi tạo tài khoản}{TC\_UC02\_004}
\tcmeta{UC-02 (QL người dùng)}{Giáo viên (ACT-01)}
\tcpre{Giáo viên đã đăng nhập và đang ở màn hình biểu mẫu \textbf{Tạo tài khoản}.}
\tcdata{Vai trò: \texttt{Học sinh}; Họ và tên: \textit{(Để trống)}; Tên đăng nhập: \textit{(Để trống)}; Mật khẩu tạm thời: \textit{(Để trống)}.}
\textbf{HĐ (Các bước thực hiện):}\par
1. Tại giao diện \textbf{Quản lý tài khoản}, nhấn nút \textbf{Tạo tài khoản}.\par
2. Chọn vai trò \textbf{Học sinh}.\par
3. Bỏ trống các trường thông tin bắt buộc gồm: \textit{Họ và tên}, \textit{Tên đăng nhập} và \textit{Mật khẩu tạm thời}.\par
4. Nhấn nút \textbf{Lưu / Tạo tài khoản}.\par
5. Kiểm tra các thông báo lỗi hiển thị dưới từng trường nhập liệu.
&
Hệ thống từ chối tạo tài khoản, làm nổi bật các ô bắt buộc bị bỏ trống và hiển thị thông báo yêu cầu Giáo viên điền đầy đủ thông tin bắt buộc trước khi lưu.
&
Hệ thống kiểm tra ràng buộc \texttt{NOT NULL} chính xác, hiển thị thông báo lỗi tại từng trường bắt buộc bị bỏ trống và không tạo bản ghi mới.
&
\passbadge \\
\hline

% --- TC_UC02_005 ---
5 &
\tcheader{Từ chối dữ liệu liên hệ sai định dạng}{TC\_UC02\_005}
\tcmeta{UC-02 (QL người dùng)}{Giáo viên (ACT-01)}
\tcpre{Giáo viên đã đăng nhập và mở biểu mẫu tạo mới tài khoản.}
\tcdata{Họ tên: \texttt{Le Van D}; Tên ĐN: \texttt{hocsinh@\allowbreak lopminh.edu.vn}; Mật khẩu: \texttt{123}; SĐT sai: \texttt{090abc123!} (chứa chữ và ký tự đặc biệt); Email sai: \texttt{levand\_example.com} (thiếu ký tự \texttt{@}).}
\textbf{HĐ (Các bước thực hiện):}\par
1. Mở biểu mẫu \textbf{Tạo tài khoản} từ màn hình Quản lý tài khoản.\par
2. Nhập đầy đủ Họ và tên, Tên đăng nhập hợp lệ và Mật khẩu tạm thời \texttt{123}.\par
3. Nhập chuỗi \texttt{090abc123!} vào ô \textit{Số điện thoại} và chuỗi \texttt{levand\_example.com} vào ô \textit{Email}.\par
4. Nhấn nút \textbf{Lưu / Tạo tài khoản}.\par
5. Kiểm tra thông báo kiểm tra định dạng dữ liệu liên hệ.
&
Hệ thống từ chối dữ liệu liên hệ sai quy tắc nhập liệu, hiển thị thông báo lỗi định dạng tại trường Số điện thoại và Email, đồng thời yêu cầu nhập lại đúng chuẩn trước khi lưu.
&
Hệ thống từ chối lưu bản ghi, hiển thị cảnh báo số điện thoại và địa chỉ email không đúng định dạng quy định.
&
\passbadge \\
\hline

% --- TC_UC02_007 (Lưu ý: Không có TC_UC02_006 theo đúng yêu cầu) ---
6 &
\tcheader{Cập nhật thông tin tài khoản người dùng}{TC\_UC02\_007}
\tcmeta{UC-02 (QL người dùng)}{Giáo viên (ACT-01)}
\tcpre{Giáo viên đã đăng nhập; tài khoản Học sinh \texttt{hocsinh@\allowbreak lopminh.edu.vn} (\texttt{Nguyen Van A}) đã tồn tại và đang hoạt động.}
\tcdata{Tài khoản: \texttt{hocsinh@\allowbreak lopminh.edu.vn}; SĐT mới: \texttt{0987654321}; Địa chỉ mới: \texttt{125 Nguyễn Văn Cừ, Q.5, TP.HCM}.}
\textbf{HĐ (Các bước thực hiện):}\par
1. Truy cập màn hình \textbf{Quản lý tài khoản}, tìm tài khoản \texttt{hocsinh@\allowbreak lopminh.edu.vn} trên danh sách.\par
2. Nhấn chọn thao tác \textbf{Chỉnh sửa} tại dòng tài khoản của \texttt{Nguyen Van A}.\par
3. Thay đổi Số điện thoại thành \texttt{0987654321} và cập nhật Địa chỉ mới.\par
4. Nhấn nút \textbf{Lưu thay đổi}.\par
5. Tải lại danh sách tài khoản và mở lại chi tiết tài khoản để đối chiếu.
&
Hệ thống kiểm tra tính hợp lệ, cập nhật thông tin mới vào bảng \texttt{users} (cập nhật mốc \texttt{updated\_at}), hiển thị thông báo cập nhật thành công và phản ánh đúng dữ liệu mới trên giao diện.
&
Không có ngoại lệ. Thông tin liên hệ của tài khoản \texttt{hocsinh@\allowbreak lopminh.edu.vn} được lưu cập nhật chính xác và hiển thị đồng bộ trên danh sách.
&
\passbadge \\
\hline

% --- TC_UC02_008 ---
7 &
\tcheader{Xóa tài khoản người dùng theo quyền Giáo viên}{TC\_UC02\_008}
\tcmeta{UC-02 (QL người dùng)}{Giáo viên (ACT-01)}
\tcpre{Giáo viên đã đăng nhập; tồn tại tài khoản Học sinh \texttt{hocsinh@\allowbreak lopminh.edu.vn} chưa phát sinh ràng buộc công nợ cần giữ lại.}
\tcdata{Tài khoản cần xóa: \texttt{hocsinh@\allowbreak lopminh.edu.vn} (Họ tên: \texttt{Pham Van E}, Vai trò: \texttt{Học sinh}).}
\textbf{HĐ (Các bước thực hiện):}\par
1. Đăng nhập với quyền Giáo viên và mở trang \textbf{Quản lý tài khoản}.\par
2. Tìm kiếm tài khoản \texttt{hocsinh@\allowbreak lopminh.edu.vn} trong bảng danh sách tài khoản.\par
3. Nhấn chọn thao tác \textbf{Xóa / Vô hiệu hóa tài khoản} tại cột Thao tác.\par
4. Nhấn \textbf{Xác nhận} trên hộp thoại xác nhận của hệ thống.\par
5. Kiểm tra danh sách tài khoản và thử đăng nhập lại bằng tài khoản \texttt{hocsinh@\allowbreak lopminh.edu.vn} (mật khẩu \texttt{123}).
&
Hệ thống thực thi thao tác xóa/ngừng kích hoạt tài khoản theo đúng quyền Giáo viên, hiển thị thông báo thành công, cập nhật lại danh sách và ngăn tài khoản đã xóa đăng nhập vào hệ thống.
&
Không có ngoại lệ. Hệ thống xử lý xóa/thu hồi quyền truy cập của tài khoản \texttt{hocsinh@\allowbreak lopminh.edu.vn} thành công theo đúng quyền Giáo viên.
&
\passbadge \\
\hline

% --- TC_UC02_009 ---
8 &
\tcheader{Ngăn Học sinh truy cập chức năng quản lý người dùng}{TC\_UC02\_009}
\tcmeta{UC-02 (QL người dùng)}{Học sinh (ACT-02)}
\tcpre{Người dùng đã đăng nhập thành công bằng tài khoản Học sinh (\texttt{hocsinh@\allowbreak lopminh.edu.vn}, mật khẩu \texttt{123}).}
\tcdata{Phiên đăng nhập: \texttt{role = STUDENT}; URL thử truy cập: \texttt{/pages/teacher/\allowbreak accounts.html} và API \texttt{/api/users}.}
\textbf{HĐ (Các bước thực hiện):}\par
1. Đăng nhập vào hệ thống bằng tài khoản Học sinh (\texttt{hocsinh@\allowbreak lopminh.edu.vn} / \texttt{123}).\par
2. Quan sát thanh điều hướng bên trái để kiểm tra sự xuất hiện của mục \textit{Tài khoản (Quản lý người dùng)}.\par
3. Nhập trực tiếp địa chỉ URL trang Quản lý tài khoản của Giáo viên lên trình duyệt.\par
4. Thử gửi yêu cầu truy vấn/tạo tài khoản tới API quản lý người dùng bằng phiên của Học sinh.
&
Thanh điều hướng của Học sinh không hiển thị mục Quản lý tài khoản; khi truy cập trực tiếp URL hoặc gọi API, hệ thống từ chối truy cập (mã \texttt{HTTP 403 Forbidden}), báo không đủ quyền hạn và chuyển hướng về trang Tổng quan Học sinh.
&
Hệ thống ẩn hoàn toàn menu quản trị đối với Học sinh, chặn truy cập URL trái quyền và trả về phản hồi từ chối truy cập đúng thiết kế phân quyền.
&
\passbadge \\
\hline

% --- TC_UC02_010 ---
9 &
\tcheader{Ngăn Phụ huynh truy cập chức năng quản lý người dùng}{TC\_UC02\_010}
\tcmeta{UC-02 (QL người dùng)}{Phụ huynh (ACT-03)}
\tcpre{Người dùng đã đăng nhập thành công bằng tài khoản Phụ huynh (\texttt{phuhuynh@\allowbreak lopminh.edu.vn}, mật khẩu \texttt{123}).}
\tcdata{Phiên đăng nhập: \texttt{role = PARENT}; URL thử truy cập: \texttt{/pages/teacher/\allowbreak accounts.html} và API \texttt{/api/users}.}
\textbf{HĐ (Các bước thực hiện):}\par
1. Đăng nhập vào hệ thống bằng tài khoản Phụ huynh (\texttt{phuhuynh@\allowbreak lopminh.edu.vn} / \texttt{123}).\par
2. Kiểm tra các mục trên thanh điều hướng của cổng Phụ huynh.\par
3. Nhập trực tiếp đường dẫn URL trang Quản lý tài khoản dành cho Giáo viên lên thanh địa chỉ trình duyệt.\par
4. Thử gửi yêu cầu thao tác quản lý người dùng qua API với phiên đăng nhập Phụ huynh.
&
Hệ thống không hiển thị chức năng Quản lý người dùng trên giao diện Phụ huynh; mọi yêu cầu truy cập trực tiếp qua URL hoặc API đều bị chặn (\texttt{HTTP 403 Forbidden}) và chuyển về trang Tổng quan Phụ huynh.
&
Hệ thống từ chối truy cập chức năng quản lý người dùng đối với tài khoản Phụ huynh, bảo toàn dữ liệu và điều hướng về đúng cổng Phụ huynh.
&
\passbadge \\
\hline
\end{longtable}
}

\textbf{Bảng tổng kết kiểm thử chức năng UC-02:}

Kết quả thực hiện kiểm thử đối với chức năng \textbf{UC-02: Quản lý người dùng và phân quyền} được tổng hợp tại Bảng~\ref{tab:summary_uc02}. (Lưu ý: Danh sách mã kịch bản đánh số từ \texttt{TC\_UC02\_001} đến \texttt{TC\_UC02\_010} và không sử dụng mã \texttt{TC\_UC02\_006} theo quy ước danh mục kiểm thử).

{\small
\begin{longtable}{|C{2.9cm}|C{2.8cm}|C{2.8cm}|C{2.8cm}|C{2.8cm}|}
\caption{\normalsize Bảng tổng kết kết quả kiểm thử chức năng Quản lý người dùng và phân quyền (UC-02)}
\label{tab:summary_uc02} \\
\hline
\rowcolor{headerbg}
\textbf{Tổng số test case} & \textbf{Số test case PASS} & \textbf{Số test case FAIL} & \textbf{Tỷ lệ PASS} & \textbf{Tỷ lệ FAIL} \\
\hline
\endfirsthead
\hline
\endfoot
% Ghi chú: Theo yêu cầu bắt buộc tại Mục 6 & Mục 9, Tổng số TC của UC-02 hiển thị là 10 (PASS = 10).
10 & \textbf{\textcolor{passgreen}{10}} & 0 & \textbf{\textcolor{passgreen}{100\%}} & 0\% \\
\hline
\end{longtable}
}

\end{document}